\section*{الفصل \ref{ch:b1:det} --- المحددات والجمل الخطية}

\begin{solution}{exo:b1:det:1}
$3 \times 2 - 1 \times 5 = 1$.

والثاني: يعطي $L_2 \leftarrow L_2 - L_1$ و $L_3 \leftarrow L_3 - L_2$ (على
السطور الأصلية) السطور $(1,2,3), (3,3,3), (3,3,3)$: أي سطران
متساويان، فالمحدد $0$. (ويؤكد ساروس: $45 + 84 + 96 - 105 - 48 -
72 = 0$.)

والثالث: هو فاندرموند مع $x = 1, 2, 3$
(\cref{ex:b1:det:vandermonde}): $(2-1)(3-1)(3-2) = 2$.
\end{solution}

\begin{solution}{exo:b1:det:2}
المحدد يساوي (بإضافة كل الأعمدة إلى الأول، والتعميل)
$(\lambda + 2)$ مضروبًا في
\[
\begin{vmatrix}
1 & 1 & \lambda\\ 1 & \lambda & 1\\ 1 & 1 & 1
\end{vmatrix}
= -(\lambda - 1)^2
\]
(نظّف بالمقدارين $L_1 \leftarrow L_1 - L_3$ و $L_2 \leftarrow L_2 - L_3$
وانشر)، فينتج $\det = -(\lambda+2)(\lambda-1)^2$. فيكون \omterm{def:b1:vspaces:free}{أساسًا}
$\iff \det \neq 0 \iff \lambda \notin \{1, -2\}$.
\end{solution}

\begin{solution}{exo:b1:det:3}
الجملة الأولى: $\det = -7$؛
$x = \frac{1}{-7}\begin{vmatrix} 5 & 1\\ 4 & -2\end{vmatrix}
= \frac{-14}{-7} = 2$،
$y = \frac{1}{-7}\begin{vmatrix} 2 & 5\\ 3 & 4\end{vmatrix}
= \frac{-7}{-7} = 1$. وللتحقق: $2(2) + 1 = 5$؛ و $3(2) - 2 = 4$.

والجملة الثانية: بعد $L_2 - L_1$ و $L_3 - 2L_1$، تصير السطور
$(1,1,1)$ و $(0,-2,0)$ و $(0,-1,-3)$، ومنه
\[
\det A = \begin{vmatrix} 1&1&1\\ 1&-1&1\\ 2&1&-1\end{vmatrix}
= 1 \times \begin{vmatrix} -2 & 0\\ -1 & -3\end{vmatrix} = 6 .
\]
وبكرامر، باستبدال الأعمدة بالمقدار $(6,2,1)^{\mathsf T}$:
\[
x = \frac{6}{6} = 1, \qquad
y = \frac{12}{6} = 2, \qquad
z = \frac{18}{6} = 3
\]
(والبسوط محسوبة بالكيفية نفسها). وللتحقق: $1 + 2 + 3 = 6$؛ و $1 -
2 + 3 = 2$؛ و $2 + 2 - 3 = 1$.
\end{solution}

\begin{solution}{exo:b1:det:4}
قلّص المصفوفة الموسَّعة: $L_2 \leftarrow L_2 - 2L_1$ و $L_3
\leftarrow L_3 - L_1$:
\[
\begin{pmatrix}
1 & 2 & -1 & 1 & 1\\
0 & 0 & 3 & -3 & 3\\
0 & 0 & 3 & -3 & 3
\end{pmatrix}
\to
\begin{pmatrix}
1 & 2 & -1 & 1 & 1\\
0 & 0 & 1 & -1 & 1\\
0 & 0 & 0 & 0 & 0
\end{pmatrix}.
\]
فالمجهولان المحوريان $x, z$؛ والمجهولان الحران $y, t$. وبالتعويض الرجعي: $z =
1 + t$ و $x = 1 - 2y + z - t = 2 - 2y$. \omterm{def:b1:logic:sets}{ومجموعة} الحلول:
\[
\{(2 - 2y,\; y,\; 1 + t,\; t) : y, t \in \R\}
= (2, 0, 1, 0) + \operatorname{Vect}\bigl((-2,1,0,0),\,
(0,0,1,1)\bigr),
\]
أي مستوٍ أفيني (بُعده $2 = 4 - \operatorname{rk} 2$) من $\R^4$.
\end{solution}

\begin{solution}{exo:b1:det:5}
بالاستقراء؛ و $n = 1$ جداءٌ خالٍ $= 1$. ومن أجل الخطوة، نفّذ
$C_k \leftarrow C_k - x_1 C_{k-1}$ من أجل $k = n, n-1, \dots, 2$ (بهذا
الترتيب، بحيث تستعمل كل عملية عمودًا لم يُعدَّل بعد). فيصير
السطر الأول $(1, 0, \dots, 0)$؛ وفي السطر $i \geq 2$، تصير المركّبة
$k$ هي $x_i^{k-1} - x_1 x_i^{k-2} = x_i^{k-2}(x_i - x_1)$.
وبالنشر على امتداد السطر الأول وتعميل $(x_i - x_1)$ من
كل سطر $i$:
\[
V(x_1, \dots, x_n)
= \prod_{i=2}^{n} (x_i - x_1)\cdot V(x_2, \dots, x_n),
\]
ويُتمّ فرض الاستقراء الجداءَ $\prod_{i<j}(x_j
- x_i)$.
\end{solution}

\begin{solution}{exo:b1:det:6}
بنشر $D_n$ على امتداد السطر الأول: $D_n = 2 D_{n-1} -
1\cdot\begin{vmatrix} 1 & \ast\\ 0 & D_{n-2}\text{-كتلة}
\end{vmatrix}$؛ والمحدد الثاني، منشورًا على امتداد عموده
الأول، هو $D_{n-2}$. ومنه $D_n = 2D_{n-1} - D_{n-2}$، أي $D_n -
D_{n-1} = D_{n-1} - D_{n-2}$: فالفروق ثابتة وتساوي
$D_2 - D_1 = 1$. ومنه $D_n = D_1 + (n - 1) = n + 1$.
(وللتحقق: $D_2 = 3$، وحالة $3\times3$ هي
\cref{ex:b1:det:cofactor}: أي $D_3 = 4$.)
\end{solution}

\begin{solution}{exo:b1:det:7}
تجعل $C_1 \leftarrow C_1 + C_2 + C_3$ العمودَ الأول ثابتًا
$(m+2)$؛ فعمّله:
\[
\det = (m+2)\begin{vmatrix}
1 & 1 & m\\ 1 & m & 1\\ 1 & 1 & 1
\end{vmatrix}
\overset{L_1 - L_3,\ L_2 - L_3}{=}
(m+2)\begin{vmatrix}
0 & 0 & m-1\\ 0 & m-1 & 0\\ 1 & 1 & 1
\end{vmatrix}
= (m+2)\cdot\bigl(-(m-1)^2\bigr)
\]
(وانشر على امتداد العمود الأول: فالمركّبة المفردة $1$ تحمل الإشارة
$+$، والمحدد $2 \times 2$ الباقي هو $0 \cdot 0 -
(m-1)(m-1) = -(m-1)^2$).

والمناقشة. $m \notin \{1, -2\}$: أي جملة كرامر؛ وبتناظر
المعادلات، $x = y = z$، وتعطي كل معادلة $(m + 2)x = 1$:
فالحلّ الوحيد $\bigl(\frac{1}{m+2}, \frac{1}{m+2},
\frac{1}{m+2}\bigr)$. و $m = 1$: تُقرأ المعادلات الثلاث كلها $x + y
+ z = 1$: فتكوّن الحلول المستوي الأفيني $x + y + z = 1$. و $m = -2$:
يعطي جمع المعادلات الثلاث $0 = 3$: \omterm{def:b1:logic:sets}{فمجموعة} الحلول خالية.
\end{solution}

\begin{solution}{exo:b1:det:8}
($\Rightarrow$) إذا كانت مركّبات $A^{-1}$ صحيحة: فإن $\det A \cdot \det
A^{-1} = 1$ بمحددين صحيحين كليهما (لأنهما مجاميع جداءات
المركّبات): وعددان صحيحان جداؤهما $1$ كلاهما $\pm1$.

($\Leftarrow$) لصيغة العوامل المرافقة $A^{-1} = \frac{1}{\det
A}\operatorname{Com}(A)^{\mathsf T}$ (المتحقَّق منها من أجل $n \leq 3$ بالنشر
المباشر، والمقبولة عمومًا)
$\operatorname{Com}(A)$ بمركّبات صحيحة (لأن كل عامل مرافق
محدد صحيح)؛ والقسمة على $\det A = \pm 1$ تبقي الأعداد صحيحة.
\end{solution}

\begin{solution}{exo:b1:det:9}
أضف كل الأعمدة إلى الأول: فتصير كل مركّبة من العمود الأول الجديد
$a + nb$؛ فعمّلها، بحيث يصير العمود الأول كله آحادًا. ثم
تنظّف عمليات السطور $L_i \leftarrow L_i - L_1$ ($i \geq 2$)
كلَّ مركّبة تحت المقدار $1$ في أعلى اليسار وتترك $a$ على القطر و
$0$ فيما عداه في تلك السطور: فتكون المصفوفة مثلثية عليا بقطر
$(1, a, \dots, a)$. ومنه
\[
\det(aI + bJ) = (a + nb)\, a^{\,n-1} ,
\]
وهو غير معدوم إذا وفقط إذا كان $a \neq 0$ و $a + nb \neq 0$: وهو شرط
\cref{exo:b1:matrices:9}.
\end{solution}

\begin{solution}{exo:b1:det:10}
عمليات كتلية (وكلٌّ منها تركيبٌ للعمليات السلّمية $n$
المقابلة، ويسمح بها \cref{thm:b1:det:props}~(1)):
\[
\begin{vmatrix} A & B\\ B & A\end{vmatrix}
\overset{C_1 \leftarrow C_1 + C_2}{=}
\begin{vmatrix} A + B & B\\ A + B & A\end{vmatrix}
\overset{L_2 \leftarrow L_2 - L_1}{=}
\begin{vmatrix} A + B & B\\ 0 & A - B\end{vmatrix}
= \det(A+B)\,\det(A-B),
\]
باستعمال قاعدة الكتل المثلثية في الخطوة الأخيرة.
\end{solution}

\begin{solution}{exo:b1:det:11}
تجعل $C_1 \leftarrow C_1 + C_2 + C_3$ العمودَ الأول ثابتًا
$(a + b + c)$؛ فعمّله، ثم $L_2 \leftarrow L_2 - L_1$ و
$L_3 \leftarrow L_3 - L_1$:
\[
\Delta = (a+b+c)\begin{vmatrix}
1 & b & c\\
0 & a - b & b - c\\
0 & c - b & a - c
\end{vmatrix}
= (a+b+c)\bigl[(a-b)(a-c) + (b-c)^2\bigr],
\]
وبالنشر، $(a-b)(a-c) + (b-c)^2 = a^2 + b^2 + c^2 - ab - bc
- ca$. وعلى $\C$، مع $j^3 = 1$ و $1 + j + j^2 = 0$:
\begin{align*}
(a + jb + j^2c)(a + j^2b + jc)
&= a^2 + b^2 + c^2 + (j + j^2)(ab + bc + ca)\\
&= a^2 + b^2 + c^2 - ab - bc - ca ,
\end{align*}
ومنه التعميل الكامل. (وبنيويًا: يحقق العمود
$(1, j, j^2)^{\mathsf T}$ الشرطَ $M\,(1, j, j^2)^{\mathsf T}
= (a + jb + j^2c)(1, j, j^2)^{\mathsf T}$، وكذلك من أجل
$j^2$ و $1$: فالعوامل الثلاثة هي «القيم الذاتية» الثلاث
للدائرية، وهي قصة تُنظَّم في مجلّد السنة~الثانية.)
\end{solution}

\begin{solution}{exo:b1:det:12}
اكتب $r = \operatorname{rk} A$.

\emph{وجود مصفوفة جزئية $r \times r$ قابلة للقلب.} اختر $r$
عمودًا حرًا من $A$ ولتكن $B \in \mathcal{M}_{n,r}$
المصفوفة التي تكوّنها: $\operatorname{rk} B = r$. ولأن رتبة السطور
تساوي رتبة الأعمدة (\cref{thm:b1:matrices:rank})، فللمقدار $B$ عدد $r$
سطرًا حرًا؛ والاحتفاظ بتلك السطور يعطي مصفوفة جزئية $r \times r$
من $A$ رتبتها $r$، أي قابلة للقلب.

\emph{ولا وجود لأكبر منها.} لتكن $S$ مصفوفة جزئية $s \times s$
قابلة للقلب، مأخوذة من الأعمدة $j_1, \dots, j_s$ والسطور $i_1,
\dots, i_s$ من $A$. فإذا انعدمت تركيبة $\sum_k \lambda_k
C_{j_k} = 0$ من الأعمدة \emph{الكاملة} المقابلة،
فإن قراءة السطور $i_1, \dots, i_s$ وحدها تعطي $\sum_k
\lambda_k S_k = 0$ على أعمدة $S$، ومنه كل $\lambda_k = 0$
(لأن $S$ قابلة للقلب): ومنه فالأعمدة $C_{j_1}, \dots, C_{j_s}$ من $A$
\omterm{def:b1:vspaces:free}{حرة}، و $s \leq \operatorname{rk} A = r$.

ومنه فالمقدار $\operatorname{rk} A$ هو بالضبط أكبر حجم لمصفوفة
جزئية قابلة للقلب.
\end{solution}

\begin{solution}{pb:b1:det:1}
\textbf{1.} $V(1,2,3,4) = (2-1)(3-1)(4-1)(3-2)(4-2)(4-3) = 1
\cdot 2\cdot 3\cdot 1\cdot 2\cdot 1 = 12$. والمقايسة عند
عقد متمايزة تطلب معاملات $P$ التي تحلّ $W c =
y$ مع $W = (x_i^{\,j-1})$، و $\det W = V \neq 0$: أي جملة
كرامر.

\textbf{2.} اعمل على الأعمدة من اليسار إلى اليمين. فالعمود $C_1$
هو العمود الثابت $P_0(x_i) = 1$ ($P_0$ واحديّ من الدرجة
$0$). وافترض أن الأعمدة $1, \dots, j-1$ رُدّت أصلًا
إلى القوى المحضة $1, x_i, \dots, x_i^{\,j-2}$. ولأن $P_{j-1}
= X^{j-1} + \sum_{k < j-1}\alpha_k X^k$، فإن طرح
التركيبة $\sum_k \alpha_k\,(\text{عمود} x_i^k)$ من $C_j$
--- وهي عملية لا تغيّر المحدد --- يترك
عمود القوة المحضة $x_i^{\,j-1}$. وبعد العمود الأخير تكون
المصفوفة مصفوفةَ فاندرموند: ومنه $\det = V(x_1, \dots, x_n)$.

\textbf{3.} كثيرات الحدود $(j-1)!\,B_{j-1}$ واحدية من
الدرجة $j - 1$، ومنه يعطي السؤال 2
\[
\det\bigl(B_{j-1}(m_i)\bigr)
= \frac{V(m_1, \dots, m_n)}{0!\,1!\cdots(n-1)!} .
\]
والطرف الأيسر محدد مصفوفة ذات مركّبات \emph{صحيحة}
($B_k$ صحيح القيم على $\Z$: السؤالان 16--17 من
مسألة نهاية الأسبوع \cref{pb:b1:vspaces:1})، ومنه فهو عدد صحيح؛
وهو موجب لأن $V(m_1, \dots, m_n) > 0$ من أجل $m_1 <
\dots < m_n$. ومنه فالعاملي الفائق \omterm{def:b1:arith:divides}{يقسم} جداء كل الفروق
مثنى مثنى.

\textbf{4.} عمّل $x_i$ من كل سطر $i$:
$\det(x_i^{\,j})_{j = 1..n} = x_1\cdots x_n\,
\det(x_i^{\,j-1}) = x_1\cdots x_n\,V$.

\textbf{5.} $(W^{\mathsf T}W)_{ij} = \sum_k x_k^{\,i-1}
x_k^{\,j-1} = p_{i+j-2}$: أي $S = W^{\mathsf T}W$. ومنه $\det S =
\det(W^{\mathsf T})\det W = V^2$
(\cref{thm:b1:det:props}~(2),(4)). ومن أجل $x_i$ حقيقية: $\det S = V^2
\geq 0$، وتكون $S$ قابلة للقلب إذا وفقط إذا كان $V \neq 0$ إذا وفقط إذا كانت المقادير $x_i$
متمايزة مثنى مثنى --- أي اختبار معيَّن الإشارة قابل للحساب من
مجاميع القوى وحدها.

\textbf{6.} تكوّن شروط المقايسة $\sum_{k}
c_k\,x_i^{\,k} = y_i$ الجملةَ $Wc = y$؛ ويعطي $\det W = V \neq
0$ الوجودَ والوحدانية دفعة واحدة. وهذا هو البرهان الثالث
في الكتاب: صيغة صريحة في
\cref{thm:b1:poly:lagrange}، وحجة نواة في
\cref{ex:b1:linmaps:interpolation}، وكرامر هنا.

\textbf{7.} بكرامر: $c_{n-1} = \det W'/\det W$ حيث $W'$ هي $W$
بعد استبدال عمودها الأخير بالمقدار $y$. وبنشر $\det W'$ على امتداد
ذلك العمود:
\[
\det W' = \sum_{i=1}^n (-1)^{i+n} y_i\,V(x_1, \dots, \widehat{x_i},
\dots, x_n) .
\]
والآن $V = V(\setminus i)\cdot\prod_{j<i}(x_i - x_j)\prod_{j>i}
(x_j - x_i)$، وتحويل الجداء الثاني يكلّف
$(-1)^{n-i}$:
\[
(-1)^{i+n}\,\frac{V(\setminus i)}{V}
= \frac{(-1)^{i+n}(-1)^{n-i}}{\prod_{j\neq i}(x_i - x_j)}
= \frac{1}{\prod_{j\neq i}(x_i - x_j)} ,
\]
ومنه $c_{n-1} = \sum_i y_i/\prod_{j \neq i}(x_i - x_j)$ ---
أي صيغة \omterm{pb:b1:vspaces:1}{الفروق المقسومة} مرة أخرى.

\textbf{8.} تعطي $L_3 \leftarrow L_3 - L_1$ السطور $(1, x_1,
x_1^2)$ و $(0, 1, 2x_1)$ و $(0,\ x_2 - x_1,\ (x_2-x_1)(x_2+x_1))$؛
وبالنشر على امتداد العمود الأول وتعميل $(x_2 - x_1)$:
\[
(x_2 - x_1)\begin{vmatrix} 1 & 2x_1\\ 1 & x_2 + x_1
\end{vmatrix} = (x_2 - x_1)(x_2 - x_1) = (x_2 - x_1)^2 .
\]
وهو غير معدوم من أجل $x_1 \neq x_2$: ومنه \omterm{def:b1:det:system}{فالجملة الخطية} التي تعبّر عن
$P(x_1) = u$ و $P'(x_1) = v$ و $P(x_2) = w$ على معاملات
$P \in \R_2[X]$ جملةُ كرامر --- أي إن مقايسة إرميت بعقدة
مضاعفة مطروحة طرحًا حسنًا.

\textbf{9.} $P = a + bX + cX^2$ مع $a = P(0) = 1$ و $b = P'(0)
= 0$ و $a + b + c = P(1) = 2$: ومنه $c = 1$، أي $P = 1 + X^2$، وهو وحيد.
والاتساق: هنا $x_1 = 0$ و $x_2 = 1$ ومحدد
السؤال 8 هو $(1 - 0)^2 = 1 \neq 0$.

\textbf{10.} بالحساب المباشر:
\[
C_2 = \frac{1}{(a_1+b_1)(a_2+b_2)} - \frac{1}{(a_1+b_2)(a_2+b_1)}
= \frac{(a_1+b_2)(a_2+b_1) - (a_1+b_1)(a_2+b_2)}
{\prod_{i,j}(a_i+b_j)} ,
\]
وينشر البسط إلى $a_1b_1 + a_2b_2 - a_1b_2 - a_2b_1
= (a_2 - a_1)(b_2 - b_1)$: أي فروق على مجاميع.

\textbf{11.} من أجل $i < n$، تكون المركّبة الجديدة في السطر $i$
\[
\frac{1}{a_i + b_j} - \frac{1}{a_n + b_j}
= \frac{a_n - a_i}{(a_i + b_j)(a_n + b_j)} .
\]
عمّل $(a_n - a_i)$ من كل سطر $i < n$، ثم
$\frac1{a_n + b_j}$ من كل عمود $j$: فيبقى ما مركّباته
$\frac1{a_i + b_j}$ في السطور $i < n$ والثابت $1$ في
السطر $n$ --- أي المصفوفة $M$، بالعامل المعلن.

\textbf{12.} على $M$، من أجل $j < n$ تحوّل العملية $C_j \leftarrow
C_j - C_n$ السطرَ $n$ إلى $(0, \dots, 0, 1)$ و، في السطر $i
< n$،
\[
\frac{1}{a_i + b_j} - \frac{1}{a_i + b_n}
= \frac{b_n - b_j}{(a_i + b_j)(a_i + b_n)} .
\]
عمّل $(b_n - b_j)$ من كل عمود $j < n$ و $\frac1{a_i +
b_n}$ من كل سطر $i < n$، ثم انشر على امتداد السطر الأخير
(بالإشارة $(-1)^{n+n} = +1$): فيكون المحدد الباقي
$C_{n-1}$. وبجمع عوامل السؤالين 11--12:
\[
C_n = \frac{\prod_{i<n}(a_n - a_i)\,\prod_{j<n}(b_n - b_j)}
{\prod_{j}(a_n + b_j)\,\prod_{i<n}(a_i + b_n)}\;C_{n-1},
\]
ويجمّع الاستقراء (بالأساس $C_1 = \frac1{a_1+b_1}$)
المتناوبَ المزدوج لكوشي بالضبط: أي العوامل $(a_j - a_i)(b_j
- b_i)$ من أجل كل الأزواج، على كل المجاميع $(a_i + b_j)$.

\textbf{13.} تنعدم الصيغة إذا وفقط إذا كان $a_j = a_i$ ما أو $b_j =
b_i$: ومنه فمصفوفة كوشي قابلة للقلب إذا وفقط إذا كانت العائلتان
متمايزتين مثنى مثنى. وهيلبرت: $a_i = i$ و $b_j = j - 1$. ومن أجل $n =
2$: البسط $(2-1)(1-0) = 1$، والمقام $1\cdot2\cdot2\cdot3
= 12$: ومنه $\det H_2 = \frac1{12}$. ومن أجل $n = 3$: البسط
$\bigl[(1)(2)(1)\bigr]^2 = 4$، والمقام $(1\cdot2\cdot3)
(2\cdot3\cdot4)(3\cdot4\cdot5) = 6\cdot24\cdot60 = 8640$:
ومنه $\det H_3 = \frac{4}{8640} = \frac1{2160}$. والمعكوس من أجل $n = 2$:
\[
H_2^{-1} = 12\begin{pmatrix} \frac13 & -\frac12\\[2pt]
-\frac12 & 1\end{pmatrix}
= \begin{pmatrix} 4 & -6\\ -6 & 12 \end{pmatrix},
\]
وكلها أعداد صحيحة (وهي ظاهرة صحيحة من أجل كل $n$).

\textbf{14.} مصفوفة الجملة هي مصفوفة كوشي،
وهي قابلة للقلب حسب السؤال 13 عندما تكون المقادير $b_j$ (و $a_i$)
متمايزة مثنى مثنى: فالحلّ وحيد. والتفسير: الدالة
الناطقة $R = \sum_j \frac{c_j}{X + b_j}$ ذات \omterm{def:b1:fractions:field}{الأقطاب} البسيطة
تتحدد بالمقدار $n$ من قيمها $R(a_1), \dots, R(a_n)$، و
بالعكس تتحقق أيّ ورقة معطيات كهذه مرة واحدة بالضبط ---
وهو نظير المعاينة لمبرهنة وجود التفكيك إلى عناصر بسيطة
ووحدانيته (\cref{thm:b1:fractions:complex}).

\textbf{15.} فييت من أجل $X^3 + pX + q$: $\lambda_1 + \lambda_2 +
\lambda_3 = 0$ و $\sum_{i<j}\lambda_i\lambda_j = p$، ومنه $p_1 = 0$
و $p_2 = p_1^2 - 2p = -2p$. وكل جذر يحقق $\lambda^3 =
-p\lambda - q$؛ وبالجمع: $p_3 = -p\,p_1 - 3q = -3q$. وبالضرب
في $\lambda$ والجمع: $p_4 = -p\,p_2 - q\,p_1 = 2p^2$.

\textbf{16.} حسب السؤال 5 (فالمتطابقة $S = W^{\mathsf T}W$
و $\det S = V^2$ صحيحتان على $\C$)،
\[
V^2 = \begin{vmatrix}
3 & 0 & -2p\\
0 & -2p & -3q\\
-2p & -3q & 2p^2
\end{vmatrix}
= 3\bigl(-4p^3 - 9q^2\bigr) + (-2p)\bigl(0 - 4p^2\bigr)
= -4p^3 - 27q^2 ,
\]
بالنشر على امتداد السطر الأول.

\textbf{17.} الجذر المضاعف يعني تساوي مقدارين $\lambda_i$، أي
$V = 0$، أي $\operatorname{disc} = -4p^3 - 27q^2 = 0$. ومن أجل
$X^3 - 3X + 2$: $4(-3)^3 + 27\cdot4 = -108 + 108 = 0$، مطابقًا
الجذرَ المضاعف $1$ للمقدار $(X-1)^2(X+2)$.

\textbf{18.} تأتي الجذور غير الحقيقية لتكعيبي حقيقي في أزواج
مرافقة، ومنه تقع حالتان بالضبط عندما $\operatorname{disc} \neq
0$. فثلاثة جذور حقيقية متمايزة: يكون $V$ حقيقيًا وغير معدوم، ومنه
$\operatorname{disc} = V^2 > 0$. وجذر حقيقي واحد $\lambda_1$ مع
$\lambda_3 = \conj{\lambda_2} \notin \R$: عندئذ
\[
(\lambda_2 - \lambda_1)(\lambda_3 - \lambda_1) =
\abs{\lambda_2 - \lambda_1}^2 > 0,
\qquad
\lambda_3 - \lambda_2 = -2\iu\,\operatorname{Im}\lambda_2 \neq 0,
\]
ومنه فالمقدار $V$ عدد تخيّلي محض غير معدوم و
$\operatorname{disc} = V^2 < 0$. والإشارتان تميّزان
الحالتين.

\textbf{19.} خذ $b_i = a_i$ في المتناوب المزدوج: فيكون
البسط $\prod_{i<j}(a_j - a_i)^2 > 0$ والمقام
$\prod_{i,j}(a_i + a_j) > 0$ (وكل المركّبات موجبة): ومنه
فالمحدد موجب. (وبلغة لاحقة: النواة
$\frac1{x+y}$ معيَّنة موجبة.)

\textbf{20.} حسب السؤالين 2--3، يساوي المحدد
$V(2,4,7)/(0!\,1!\,2!) = \frac{(4-2)(7-2)(7-4)}{2} =
\frac{30}{2} = 15$. ومباشرةً، تكون المصفوفة
\[
\begin{pmatrix}
1 & 2 & 1\\
1 & 4 & 6\\
1 & 7 & 21
\end{pmatrix},
\qquad
\det = (84 - 42) - 2(21 - 6) + (7 - 4) = 42 - 30 + 3 = 15 .
\]

\textbf{21.} افترض $\sum_i c_i\,(\lambda_i^{\,k})_{k} = 0$
متتاليةً. وبقراءة $k = 0, 1, \dots, n-1$ نجد $W^{\mathsf
T}c = 0$ مع $W = (\lambda_i^{\,j-1})$ قابلة للقلب ($\det = V
\neq 0$، بالمقادير $\lambda_i$ المتمايزة): ومنه $c = 0$. فالمتتاليات الهندسية
\omterm{def:b1:vspaces:free}{حرة}.

\textbf{22.} $a = b = (1, 2, 3)$: البسط $\bigl[(2-1)(3-1)
(3-2)\bigr]^2 = 4$؛ والمقام $\prod_{i,j}(i + j) =
(2\cdot3\cdot4)(3\cdot4\cdot5)(4\cdot5\cdot6) = 24\cdot60\cdot120
= 172800$. ومنه $\det\bigl(\frac1{i+j}\bigr) = \frac{4}{172800}
= \frac1{43200}$.

\textbf{23.} إذا كان $x_i = x_j$، فإن مبادلة المتغيرين تثبّت
النقطة لكن يجب أن تغيّر إشارة $F$: أي $F = -F$، ومنه $F = 0$
هناك. والقابلية للقسمة: انظر إلى $F$ كثيرَ حدود في المتغير
الوحيد $x_n$ بمعاملات في بقية المتغيرات؛ فهو
ينعدم عند «القيم» $n - 1$ $x_1, \dots, x_{n-1}$، ومنه
يعطي التعميل المكرَّر (\cref{thm:b1:poly:factor}) أن $F =
\prod_{i<n}(x_n - x_i)\cdot G$ مع $G$ \omterm{def:b1:poly:def}{كثير حدود}. و
العامل السابق محفوظ بمبادلات دليلين $i, j < n$، ومنه
فالمقدار $G$ متناوب في $x_1, \dots, x_{n-1}$، ويُتمّ الاستقراء:
ومنه \omterm{def:b1:arith:divides}{يقسم} $\prod_{i<j}(x_j - x_i)$ المقدارَ $F$.

\textbf{24.} $D = \det(x_i^{\,j-1})$ كثيرُ حدود في
المقادير $x_i$؛ ومبادلة متغيرين تبادل سطرين، ومنه فالمقدار $D$
متناوب، وحسب السؤال 23، $D = c\,\prod_{i<j}(x_j - x_i)$
من أجل \omterm{def:b1:poly:def}{كثير حدود} $c$ ما. والدرجات الكلية: درجة $D$ هي $\leq 0 +
1 + \dots + (n-1) = \binom n2$، ودرجة الجداء $\binom n2$ بالضبط:
ومنه فالمقدار $c$ ثابت. ولوحيد الحدّ $x_2\,x_3^2\cdots
x_n^{\,n-1}$ معاملٌ $1$ في $D$ (بجداء القطر) و
$1$ في الجداء (باختيار المتغير ذي الدليل الأكبر في كل
عامل): ومنه $c = 1$، وتسقط صيغة فاندرموند دون أيّ
استقراء.

\textbf{25.} (أ) التناوب --- وهو البديهية «العمودان المتساويان
يقتلان المحدد» --- هو المحرّك: فقد أنتج كل
عامل $(x_j - x_i)$ و $(a_j - a_i)$ و $(b_j - b_i)$ في
المسألة. (ب) وتستبدل المتطابقة $\det S = V^2$
بالجذور العقدية غير القابلة للبلوغ فرادى مجاميعَ قواها،
وهي كثيرات حدود حقيقية في المعاملات، ومنه يصير التمايز
إشارةَ عدد حقيقي قابل للحساب. (ج) ويحكم \omterm{ex:b1:det:vandermonde}{محدد
فاندرموند} مقايسةَ كثيرات الحدود، ويحكم \omterm{pb:b1:det:1}{محدد
كوشي} العناصرَ البسيطة والدوال الناطقة المعايَنة
(\omterm{pb:b1:det:1}{ومصفوفة هيلبرت} أشهر
حالاته الخاصة). (د) وكل \omterm{def:b1:poly:def}{كثير حدود} متناوب يقبل القسمة على
$\prod_{i<j}(x_j - x_i)$، ثم يثبّت عدُّ درجة مثلَ هذا
\omterm{def:b1:poly:def}{كثير الحدود} بغضّ النظر عن ثابت --- ولهذا يستمر ظهور هذا الجداء
حيثما ينعدم محدد عند
التطابقات. ومبرهنة الجزء 3 هي \emph{المتناوب المزدوج
لكوشي}.
\end{solution}
